\documentclass[11pt]{article}
\usepackage[margin=1in]{geometry}
\usepackage{amsmath,amssymb,amsthm}
\usepackage[colorlinks=true,urlcolor=blue,linkcolor=blue,citecolor=blue]{hyperref}
\newtheorem{theorem}{Theorem}
\newtheorem{lemma}{Lemma}
\title{A log-convexity proof of the\\converse binomial inequality}
\author{Principia Math in Progress}
\date{12 September 2026}
\begin{document}
\maketitle

\paragraph{Task 1: Resolve MathDB 352508.}
We prove Gerhold--Simon's Conjecture 3.2 \cite{GS}, with the sharp
constant, all equality cases, and a stronger bound depending on \(x/y\).
As in the source, generalized binomial coefficients are defined by gamma
functions.

\begin{theorem}\label{thm:main}
For every positive integer \(k\), every real \(\alpha>2\), and \(x,y\geq0\),
\begin{equation}\label{eq:main}
 2^{\alpha-1}\sum_{j=0}^k
 \frac{\Gamma(\alpha k+1)}
      {\Gamma(\alpha j+1)\Gamma(\alpha(k-j)+1)}
 x^{\alpha j}y^{\alpha(k-j)}
 \ \geq\ (x+y)^{\alpha k}.
\end{equation}
The constant \(2^{\alpha-1}\) is best possible uniformly in \(k,x,y\).
Equality holds precisely when \(x=y=0\), or when \(k=1\) and \(x=y>0\).

More precisely, if \(x+y>0\) and \(p=x/(x+y)\), then
\begin{equation}\label{eq:refined}
 \sum_{j=0}^k\binom{\alpha k}{\alpha j}
       p^{\alpha j}(1-p)^{\alpha(k-j)}
 \ \geq\
 \left(\frac{1+(2p-1)^{2k}}2\right)^{\alpha-1}.
\end{equation}
For \(k\geq2\) and \(0<p<1\), this refinement is strict.
\end{theorem}

\paragraph{The gamma quotient is log-convex along a ray.}
\begin{lemma}\label{lem:gamma}
If \(A,B>0\), then
\[
 t\longmapsto
 \frac{\Gamma(1+(A+B)t)}{\Gamma(1+At)\Gamma(1+Bt)}
\]
is strictly log-convex for \(t>0\).
\end{lemma}
\begin{proof}
The standard trigamma series, or its termwise Laplace integral, gives
\[
 \psi'(1+z)=\sum_{n=1}^{\infty}\frac1{(n+z)^2}
          =\int_0^\infty \frac{u e^{-zu}}{e^u-1}\,du,
 \qquad z>0.
\]
After rescaling the three integrals, the second derivative of the
logarithm of the quotient is
\begin{equation}\label{eq:kernel}
 \int_0^\infty u e^{-tu}
 \left\{
 \frac1{e^{u/(A+B)}-1}
 -\frac1{e^{u/A}-1}
 -\frac1{e^{u/B}-1}
 \right\}\,du.
\end{equation}
The expression in braces is strictly positive.
Indeed, for fixed \(u>0\), put
\[
 f_u(c)=\frac1{e^{u/c}-1},\qquad c>0.
\]
Then \(f_u(c)/c=u^{-1}v/(e^v-1)\), where \(v=u/c\).
The function \(v/(e^v-1)\) is strictly decreasing on \(v>0\):
the numerator of its derivative is
\((1-v)e^v-1<0\), since this numerator vanishes at zero and
has derivative \(-ve^v<0\).
Thus \(f_u(c)/c\) increases strictly in \(c\), and
\[
 f_u(A)+f_u(B)
 =A\frac{f_u(A)}A+B\frac{f_u(B)}B
 <(A+B)\frac{f_u(A+B)}{A+B}=f_u(A+B).
\]
This proves positivity in \eqref{eq:kernel}. All integrals converge:
the expression in braces tends to \(1/2\) at zero and decays
exponentially at infinity. The second logarithmic derivative is
therefore strictly positive.
\end{proof}

\paragraph{Log-convexity of the finite binomial sum.}
Fix \(k\geq1\) and \(p\in(0,1)\), and define
\[
 S(t)=\sum_{j=0}^k T_j(t),\qquad
 T_j(t)=\binom{tk}{tj}p^{tj}(1-p)^{t(k-j)},\qquad t>0.
\]
For \(1\leq j\leq k-1\), Lemma \ref{lem:gamma} with
\(A=j\), \(B=k-j\) proves strict log-convexity of \(T_j\).
The endpoint terms are exponentials in \(t\), and hence log-convex.
The positive factors involving \(p\) add only affine functions to
the logarithms and do not change their second derivatives.

A positive finite sum of log-convex functions is log-convex. To check
both this assertion and the strictness needed here, let
\(\ell_j=\log T_j\) and \(\omega_j=T_j/S>0\). Direct differentiation gives
\begin{equation}\label{eq:sum}
 (\log S)''
 =\sum_{j=0}^k\omega_j\ell_j''
  +\sum_{j=0}^k\omega_j
       \left(\ell_j'-\sum_{i=0}^k\omega_i\ell_i'\right)^2\geq0.
\end{equation}
When \(k\geq2\), an interior summand makes the first term strictly
positive, so \(\log S\) is strictly convex on \((0,\infty)\).

\paragraph{Use the ordinary binomial identities at \(t=1,2\).}
The usual binomial theorem and its even-part identity give
\begin{equation}\label{eq:anchors}
 S(1)=1,\qquad
 S(2)=\frac{1+(2p-1)^{2k}}2\geq\frac12.
\end{equation}
For \(\alpha>2\), write
\[
 2=\frac{\alpha-2}{\alpha-1}\,1+
                       \frac1{\alpha-1}\,\alpha.
\]
Convexity of \(\log S\) implies
\[
 \log S(2)\leq
 \frac{\alpha-2}{\alpha-1}\log S(1)
 +\frac1{\alpha-1}\log S(\alpha)
 =\frac1{\alpha-1}\log S(\alpha).
\]
Consequently
\[
 S(\alpha)\geq S(2)^{\alpha-1}
 =\left(\frac{1+(2p-1)^{2k}}2\right)^{\alpha-1}
 \geq2^{1-\alpha}.
\]
This proves \eqref{eq:refined}, and multiplication by
\(2^{\alpha-1}(x+y)^{\alpha k}\) proves \eqref{eq:main}.
For \(p=0\) or \(p=1\), \(S(t)=1\), so the refinement holds directly
and the original inequality is strict when \(x+y>0\).

\paragraph{Strictness and sharpness.}
If \(k\geq2\) and \(0<p<1\), the convexity step is strict by
\eqref{eq:sum}. If \(k=1\) and \(p\ne1/2\), then \(S(2)>1/2\),
which already makes the original inequality strict.
If \(k=1\) and \(p=1/2\), then \(S(\alpha)=2^{1-\alpha}\), giving
equality and proving sharpness of the uniform constant.
Finally, \(x=y=0\) gives equality by continuity, since all terms
have positive total degree. This accounts for every case.

\paragraph{Verification and status.}
The companion verifier checks ordinary binomial identities and
integer-parameter refinements with exact arithmetic, and independently
checks the noninteger sums, trigamma derivatives, and integral kernel.
These computations supplement the proof.
The current MathDB entry, the current source, its published HTML
Conjecture 3.2, the author's addenda, and subsequent public claims were
checked. No matching full solution claim was located; the dated audit
records the unrelated catalogue hits and later fractional-binomial work.

\paragraph{Task 2: Determine the sharp constants after fixing \(k\).}
For each \(k\geq2\) and \(\alpha>2\), determine
\(\min_{0\leq p\leq1}S(\alpha)\), its minimizing points, and the
corresponding best constant in \eqref{eq:main}.

\paragraph{Task 3: Find optimal multinomial converse constants.}
Replace the two variables by an arbitrary number of nonnegative
variables and determine the smallest uniform constant multiplying
the generalized multinomial sum.

\paragraph{Task 4: Classify shifted gamma quotients.}
Determine which positive shifts preserve log-convexity of the quotient
in Lemma \ref{lem:gamma}, and use those ranges to establish converse
inequalities for sums with shifted generalized binomial coefficients.

\paragraph{Task 5: Quantify the strict interpolation gap.}
Derive explicit positive lower bounds for
\(\log S(\alpha)-(\alpha-1)\log S(2)\) when \(k\geq2\),
with useful control near \(\alpha=2\) and near \(p=0,1\).

\begin{thebibliography}{9}
\bibitem{GS}
S.\ Gerhold and T.\ Simon,
A converse to the neo-classical inequality with an application to the
Mittag-Leffler function,
\emph{Monatshefte f\"ur Mathematik} \textbf{200} (2023), 627--645.
\href{https://doi.org/10.1007/s00605-022-01817-8}{doi:10.1007/s00605-022-01817-8}.
\href{https://arxiv.org/abs/2111.02747}{arXiv:2111.02747v2},
Conjecture 3.2.
\end{thebibliography}
\end{document}
